\documentclass[conference]{IEEEtran}

\usepackage{amsmath,amssymb,amsfonts}
\usepackage{algorithmic}
\usepackage{algorithm}
\usepackage{array}
\usepackage{booktabs}
\usepackage{tabularx}
\usepackage{listings}
\usepackage{xcolor}
\usepackage{url}
\usepackage{cite}
\usepackage{microtype}

\newcolumntype{Y}{>{\centering\arraybackslash}X}
\definecolor{codebg}{rgb}{0.97,0.97,0.97}
\definecolor{codeborder}{rgb}{0.85,0.85,0.85}
\definecolor{codecomment}{rgb}{0.20,0.50,0.20}
\definecolor{codekeyword}{rgb}{0.12,0.25,0.70}
\definecolor{codestring}{rgb}{0.65,0.15,0.15}

\lstdefinestyle{conferenceListing}{
  backgroundcolor=\color{codebg}, frame=single,
  rulecolor=\color{codeborder}, framerule=0.5pt, framesep=2.5pt,
  commentstyle=\color{codecomment}\itshape,
  keywordstyle=\color{codekeyword}\bfseries,
  stringstyle=\color{codestring}, basicstyle=\ttfamily\scriptsize,
  breaklines=true, breakatwhitespace=false, captionpos=b,
  numbers=left, numbersep=4pt, numberstyle=\tiny\color{gray},
  tabsize=2, showstringspaces=false, xleftmargin=1.0em, xrightmargin=0.2em
}
\lstset{style=conferenceListing}
\setlength{\tabcolsep}{2.8pt}
\renewcommand{\arraystretch}{1.15}

\begin{document}

\title{\textbf{\Large CloudGuard: Explainable Policy Violation Detection and Remediation Guidance for AWS Infrastructure as Code}}

\author{
  \IEEEauthorblockN{\textbf{Uday and Project Team}}
  \IEEEauthorblockA{\textit{Department of Computer Science and Engineering}\\
  \textit{Mohan Babu University}\\ Tirupati, India}
}

\maketitle

\begin{abstract}
Infrastructure as Code (IaC) makes cloud deployment repeatable, but insecure configuration can be propagated at the same speed. This paper presents CloudGuard, a browser-based prototype for local analysis of AWS Terraform and CloudFormation files. CloudGuard detects five policy conditions: public S3 access, publicly accessible RDS, unrestricted security-group ingress, wildcard IAM actions, and disabled RDS encryption. Each finding includes a policy identifier, severity, risk deduction, control reference, and remediation guidance. The implementation is intentionally scoped to guidance and transparent scoring; it does not claim automatic repair, Kubernetes or Helm support, expert annotation, or human-subject evidence. A reproducible project-team synthetic mutation benchmark contains 50 manifests and 70 labelled policy-instance pairs. On this benchmark, the prototype produced 30 true positives, 0 false positives, 0 false negatives, and 40 true negatives, yielding precision, recall, and F1 of 1.0000. These values are reported only as an implementation sanity check because the benchmark was constructed around the same five rules.
\end{abstract}

\begin{IEEEkeywords}
Infrastructure as Code, cloud security, policy violation detection, explainable risk, Terraform, CloudFormation.
\end{IEEEkeywords}

\section{Introduction}
Cloud infrastructure is increasingly declared in machine-readable configuration rather than provisioned manually. This improves repeatability, but it also allows a permissive access rule or disabled protection to be copied across environments before deployment review. Security scanners can identify many conditions, yet their output is often difficult for a developer or reviewer to connect to a concrete policy and next action.

CloudGuard addresses this traceability problem with a small, inspectable AWS IaC analyzer. Security scanners such as Checkov, KICS, and Trivy are commonly used for IaC security analysis~\cite{checkov,kics,trivy}. The prototype reads Terraform and CloudFormation files locally in the browser, applies five explicit rules, computes a transparent score from configured deductions, and displays remediation guidance. It is a capstone research artifact and demonstration, not a production cloud security service.

Minna et al.~\cite{minna2026automated} investigated inconsistent security-policy analysis and policy mapping across analyzers for Helm charts in Kubernetes environments. CloudGuard studies the related traceability problem for AWS Terraform and CloudFormation, without implementing Helm or Kubernetes support.

The contributions are: (1) a browser workflow for local AWS IaC policy review; (2) a finding representation that preserves policy ID, severity, risk points, control reference, and remediation guidance; (3) a reproducible synthetic mutation benchmark for implementation sanity testing; and (4) an explicit boundary between measured results and future experiments.

\section{Scope and Research Questions}
CloudGuard supports Terraform and CloudFormation~\cite{terraform,cloudformation} only. It does not support Helm or Kubernetes, does not execute changes against AWS, and does not automatically rewrite source files. The selected file is read in the browser, matched against deterministic rules, and rendered as a report and remediation view.

We define three questions for the current prototype evaluation:
\begin{itemize}
  \item \textbf{RQ1:} Does the rule set identify the policy mutations represented in a frozen synthetic AWS IaC benchmark?
  \item \textbf{RQ2:} Can findings be presented with explainable severity, risk points, and remediation guidance?
  \item \textbf{RQ3:} Can controlled remediation guidance remove the targeted policy violations without introducing additional structural errors in the synthetic benchmark?
\end{itemize}

\section{System Design}
The user selects a \texttt{.tf}, \texttt{.yaml}, \texttt{.yml}, or \texttt{.json} file, or loads the built-in vulnerable sample. The browser places the text in the editor. Matching rules become finding records. The renderer displays a report, a 0--100 score, and remediation recommendations. No cloud credentials or deployment operation is required.

The current policies are S3-001 for public S3 access, RDS-001 for public RDS access, SG-001 for unrestricted ingress, IAM-001 for wildcard IAM actions, and RDS-002 for disabled database encryption. The selected controls are aligned with common cloud security guidance, including the AWS Well-Architected Security Pillar~\cite{awswell}. The score begins at 100 and subtracts 28, 25, 18, 18, and 14 points for these findings, respectively. These deductions are CloudGuard-defined heuristic weights selected to provide transparent relative prioritization; they are not calibrated probabilities of compromise.

\begin{table}[htbp]
\centering
\caption{Implemented CloudGuard Policies and Deductions}
\label{tab:policies}
\scriptsize
\begin{tabularx}{\columnwidth}{lXcc}
\toprule
\textbf{Policy ID} & \textbf{Condition} & \textbf{Severity} & \textbf{Points}\\
\midrule
S3-001 & Public S3 access & Critical & 28\\
RDS-001 & Public RDS access & Critical & 25\\
SG-001 & Unrestricted ingress & High & 18\\
IAM-001 & Wildcard IAM actions & High & 18\\
RDS-002 & Disabled encryption & High & 14\\
\bottomrule
\end{tabularx}
\end{table}

\section{Benchmark and Reproducibility}
The research package contains a 50-manifest project-team synthetic mutation benchmark: 25 Terraform and 25 CloudFormation manifests, with 70 labelled policy-instance pairs. A policy instance associates a file, resource, format, policy SID, expected violation label, and severity. This avoids treating one file as one result when several policies occur in the same file.

The labels were produced by the project team using predefined criteria. Two reviewer files are retained as an audit trail, but the collection process is not documented as independent blinded annotation. Therefore, this paper does not report Cohen's kappa or claim inter-rater reliability. The benchmark is suitable for regression testing and implementation sanity checks, but not for generalization to arbitrary production configurations.

Results are regenerated from retained labels and predictions. The public repository is \url{https://github.com/duday8329-ai/cloudguard-research}.

\section{Evaluation Method}
For CloudGuard, each policy-instance prediction is compared with the expected violation label. Precision, recall, and F1 are calculated as
\begin{equation}
\mathrm{Precision}=\frac{TP}{TP+FP}, \qquad
\mathrm{Recall}=\frac{TP}{TP+FN}.
\end{equation}
\begin{equation}
F1=2\frac{\mathrm{Precision}\times\mathrm{Recall}}
{\mathrm{Precision}+\mathrm{Recall}}.
\end{equation}
Baseline-analyzer metrics are reserved for future comparison after validated semantic mappings are established. Unmapped findings are not silently treated as negatives. Latency is measured over repeated local analyzer runs and reported with the environment and trial count.

\begin{lstlisting}[language=bash,caption={Example insecure Terraform ingress configuration.},label={lst:insecure}]
resource "aws_security_group" "bastion_ingress" {
  ingress {
    from_port   = 22
    to_port     = 22
    protocol    = "tcp"
    cidr_blocks = ["0.0.0.0/0"]
  }
}
\end{lstlisting}

\section{Results}
The CloudGuard sanity run generated 30 true positives, no false positives, no false negatives, and 40 true negatives over 70 policy-instance labels. Precision, recall, and F1 are all 1.0000. Because the mutation cases were designed to exercise the same five prototype rules, this result is expected to be optimistic. It demonstrates that the pipeline, labels, predictions, and metric generation work together; it does not demonstrate perfect performance on unseen AWS IaC.

\begin{table}[htbp]
\centering
\caption{Reproducible Results Currently Available}
\label{tab:results}
\scriptsize
\begin{tabularx}{\columnwidth}{lYY}
\toprule
\textbf{Metric} & \textbf{Measured Value} & \textbf{Interpretation}\\
\midrule
Manifests & 50 & 25 Terraform + 25 CloudFormation synthetic mutations\\
Policy instances & 70 & Project-team labels at policy-instance level\\
TP / FP / FN / TN & 30 / 0 / 0 / 40 & Generated CloudGuard sanity matrix\\
Precision / Recall / F1 & 1.0000 / 1.0000 / 1.0000 & Optimistic synthetic result; not generalization evidence\\
Latency & 10 trials; median 66.091 ms; p95 84.382 ms & Local Windows prototype measurement\\
\bottomrule
\end{tabularx}
\end{table}

For RQ2, the prototype presents each finding with a stable policy ID, severity, risk deduction, control reference, and remediation guidance, demonstrating the intended explainability workflow. Raw baseline scans for Checkov, KICS, and Trivy are retained, but this paper makes no broad superiority claim. Comparative metrics require validated semantic mappings for the compared policy instances.

\section{Controlled Remediation Guidance Validation}
Thirty positive policy instances were copied into before/after cases. For each target case, a manually prepared remediation was applied to the corresponding configuration, after which CloudGuard re-scanned the modified configuration: public access was restricted, encryption was enabled, unrestricted CIDRs were narrowed, and wildcard actions were replaced with a specific action. For RQ3, manually prepared remediation guidance was successfully applied to all 30 controlled cases, and CloudGuard detected no remaining target violations or recorded structural errors. Native Terraform and CloudFormation validators were unavailable, so Gate 1 is a structural integrity check. The result measures controlled remediation guidance clearance and does not establish automatic repair capability.

\section{Leave-One-Policy-Out Ablation}
The full system produced precision 1.0000, recall 1.0000, and F1 1.0000. Removing any one of the five policy rules produced precision 1.0000, recall 0.8000, and F1 0.8889, with six false negatives. Each policy family contributes six positive cases in this balanced synthetic benchmark. These results show sensitivity to removal of individual rules within the synthetic benchmark, not general model robustness.

\section{Ranking Consistency Check}
The CloudGuard-ranked findings produced NDCG@10 of 1.0000. The relevance labels were assigned with the predefined severity and risk-point rubric: public S3 and public RDS received 4, unrestricted ingress and wildcard IAM received 3, and disabled encryption received 2. This is a consistency check against the same rubric used by the ranking, not evidence of independent ranking quality or expert-consensus NDCG.

\section{Limitations and Evidence Boundary}
The current package does not contain a human-participant study, production configurations, automatic AST-based repair, native IaC validation logs, or a large external benchmark. The implementation is pattern-based and can miss equivalent configurations expressed through variables, modules, intrinsic functions, indirection, or uncommon syntax. It does not build a full Terraform or CloudFormation semantic model. Remediation output is advice rather than an AST transformation, and the score is a transparent prioritization aid rather than a calibrated risk model.

\section{Future Evaluation Protocol}
A larger study should add manually reviewed Terraform and CloudFormation examples while preserving policy-instance ground truth. The same files should be scanned by CloudGuard, Checkov, KICS, and Trivy using a published normalized mapping. Ranking relevance should be assigned before any system ranking is shown. Remediation cases should use native Terraform and CloudFormation validators. Any developer study requires consent and appropriate oversight.

\section{Conclusion}
CloudGuard is a compact, explainable AWS IaC analyzer ready for demonstration and reproducible synthetic evaluation. Its current strengths are transparency, stable policy IDs, visible deductions, local file analysis, generated metrics, controlled remediation guidance checks, and component-sensitivity results. The benchmark supports an honest prototype evaluation, not a broad production-accuracy claim. The ranking value is a consistency check against the predefined rubric, not independent ranking-quality evidence. These results support CloudGuard as a reproducible prototype for AWS IaC policy analysis while leaving broader comparative and deployment claims for future evaluation.

\begin{thebibliography}{7}
\bibitem{minna2026automated} F.~Minna, A.~Blaise, K.~Tuma, and F.~Massacci, ``Automated analysis of security policy violations in Helm charts,'' \emph{IEEE Trans. Dependable and Secure Computing}, vol.~23, no.~2, pp. 2519--2533, 2026, doi: 10.1109/TDSC.2025.3628213.
\bibitem{awswell} Amazon Web Services, ``AWS Well-Architected Framework: Security Pillar,'' \url{https://docs.aws.amazon.com/wellarchitected/latest/security-pillar/welcome.html}.
\bibitem{checkov} Bridgecrew, ``Checkov Documentation,'' \url{https://www.checkov.io/}.
\bibitem{kics} Checkmarx, ``KICS Documentation,'' \url{https://docs.kics.io/}.
\bibitem{trivy} Aqua Security, ``Trivy Documentation,'' \url{https://trivy.dev/}.
\bibitem{terraform} HashiCorp, ``Terraform Language Documentation,'' \url{https://developer.hashicorp.com/terraform/language}.
\bibitem{cloudformation} Amazon Web Services, ``AWS CloudFormation User Guide,'' \url{https://docs.aws.amazon.com/AWSCloudFormation/latest/UserGuide/Welcome.html}.
\end{thebibliography}

\end{document}
